\documentclass[11pt]{article}
\usepackage[margin=1in]{geometry}
\usepackage{amsmath,amssymb,amsthm,mathtools,booktabs,array,enumitem}
\usepackage[T1]{fontenc}
\usepackage{lmodern,microtype}
\usepackage[colorlinks=true,linkcolor=blue,urlcolor=blue,citecolor=blue]{hyperref}
\allowdisplaybreaks
\setlist{itemsep=3pt,topsep=5pt}
\newtheorem{proposition}{Proposition}
\newtheorem{theorem}[proposition]{Theorem}
\newtheorem{corollary}[proposition]{Corollary}
\newcommand{\ii}{\mathrm{i}}
\newcommand{\bbone}{\boldsymbol{1}}
\newcommand{\btheta}{\boldsymbol{\theta}}
\newcommand{\bk}{\boldsymbol{k}}
\newcommand{\bh}{\boldsymbol{h}}
\newcommand{\bx}{\boldsymbol{x}}
\newcommand{\by}{\boldsymbol{y}}
\newcommand{\bDelta}{\boldsymbol{\Delta}}
\newcommand{\cH}{\mathcal H}
\title{Using the Surplus Points in a Lattice Sequence\\
\large A proposed direction for revision and discussion}
\author{Working note for the authors}
\date{October 24, 2026}
\begin{document}
\maketitle
\begin{abstract}
The most promising revision combines referee Z's suggestions 2 and 3:
make the use of surplus points computationally practical, prove a precise
improvement over truncation, and test that improvement for smoothness
$\alpha>1$. Megan Millet's proposed weights, constant on each binary block,
are a natural starting point. However, that class, its power weights, and
its optimal weights already appear in Hickernell--Kritzer--Kuo--Nuyens
(2012; HKKN12). This note derives a specialized binary covariance
identity and a scalar recursion for the same optimal block rule.
For a positive-spectrum periodic Korobov kernel, the recursion gives
strictly positive weights at every available point and a strictly smaller
Hilbert-space worst-case error than truncation at every non-dyadic
sample size. This is a candidate contribution, not a verified claim of
novelty. It does not remove every truncation constant or establish the
analogous Banach-space improvement. Small exact examples clarify
those limitations.
\end{abstract}

\section{The revision should answer a different main question}
Referees Y and Z identify a substantive issue: an arbitrary-$n$ rate
does not by itself justify using all $n$ values.
Put $M=\lfloor\log_2n\rfloor$ and $N=2^M$.
Keeping only the first $N$ values already gives the known dyadic bound,
while discarding fewer than half the evaluations.
Full nodewise optimized weights can reproduce that rule with zero
weights on the surplus points. Their non-increasing optimal error is
therefore an immediate feasibility fact.

A more focused question is:
\begin{quote}
Can all completed values be given positive weights, with inexpensive
postprocessing, and can their error reduction over truncation
be quantified and demonstrated for smooth functions?
\end{quote}
I recommend centering the revision on the periodic Hilbert-space
quantity $P_{\alpha,2}$, retaining corrected equal-weight Banach analysis
as context. Keep the two spaces and their smoothness parameters explicit:
an improvement in one is not a theorem about the other.

\paragraph{What is already in HKKN12.}
Sections 3.1, 3.3, 4.2 and 4.3 of \cite{HKKN12} give optimal compound
weights, explicit power weights, the base-$b$ block partition, and optimal
compound lattice rules. Proposition 2, pp.~174--176, already reduces a
block covariance from a double sum to a sum over the larger block.
Neither $13=8+4+1$ nor solving the small quadratic optimization problem
is a new construction. The possible additional contribution below is
the binary specialization, a positive-weight scalar recursion, an exact
gain over truncation, and a reproducible computational study.
Further literature checking is needed before calling even that
specialization new.

\section{Set-up and the binary covariance identity}
For this note, $\alpha>1$ and $\gamma_j>0$.
Consider the manuscript's periodic Korobov Hilbert space, with
\[
 K(\bx,\by)=\kappa(\bx-\by),\qquad
 \kappa(\bx)=1+\sum_{\bk\ne0}\rho_{\bk}
                e^{2\pi\ii\bk^\top\bx},\qquad
 \rho_{\bk}=\widetilde r(\bk,\boldsymbol{\gamma})^{-2\alpha}>0.
\]
Here $\widetilde r(\bk,\boldsymbol{\gamma})=\prod_{j:k_j\ne0}(|k_j|/\gamma_j)$,
with empty product one. The sum is absolutely convergent and $\int K(\bx,\by)\,d\bx=1$.
Let $e(Q)=\sup_{\|f\|_{\cH}\le1}|Q(f)-I(f)|$.
For node weights summing to one,
\begin{equation}\label{eq:fourier}
 e(Q)^2=\sum_{\bk\ne0}\rho_{\bk}
 \left|\sum_iw_i e^{2\pi\ii\bk^\top\bx_i}\right|^2.
\end{equation}
A common shift $\bDelta$ does not affect this error.

Let $\bx_i=\{\bh\phi_2(i)+\bDelta\}$ be a binary rank-1 lattice sequence,
with an odd coordinate of $\bh$, say $h_1$.
On the finite horizon through $2^{M+1}$ points, take a coherent integer
representative of the generator modulo $2^{M+1}$.
An existence theorem for a good infinite generator is not a guarantee
for every software default.
Define
\[
 H_m=\{\bh j/2^m\bmod1:0\le j<2^m\},\qquad
 \mathcal D_m=\{\bk\ne0:\bk^\top\bh=0\bmod2^m\}.
\]
Then $|H_m|=2^m$ and $H_m\subset H_{m+1}$.
For the full equal-weight rule, write
\begin{equation}\label{eq:dm}
 d_m=e(Q_{2^m})^2
 =2^{-m}\sum_{j=0}^{2^m-1}\kappa(\bh j/2^m)-1
 =\sum_{\bk\in\mathcal D_m}\rho_{\bk}.
\end{equation}
Here $d_m>d_{m+1}>0$. For strict decrease, the positive coefficient
at $\bk=2^m\boldsymbol e_1$ is in
$\mathcal D_m\setminus\mathcal D_{m+1}$.

Write the nonzero binary digits of $n$ as
\[
 n=N_1+\cdots+N_L,\quad N_r=2^{m_r},\quad
 m_1>\cdots>m_L\ge0,\quad a_r=\sum_{q<r}N_q.
\]
Set $B_r=\{a_r,\ldots,a_r+N_r-1\}$ and
$Q_r(f)=N_r^{-1}\sum_{i\in B_r}f(\bx_i)$.
Megan's block-constant rule is
\begin{equation}\label{eq:compound}
 Q_{\btheta}(f)=\sum_{r=1}^L\theta_rQ_r(f),\qquad
 \sum_r\theta_r=1,\qquad w_i=\theta_r/N_r\quad(i\in B_r).
\end{equation}
The $\theta_r$ are \emph{block masses}, not individual sample weights.
Equal point weights have $\theta_r=N_r/n$;
HKKN12 power weights have $\theta_r=N_r^a/\sum_qN_q^a$;
truncation has $\btheta=(1,0,\ldots,0)$.

\begin{proposition}[Binary block covariance]\label{prop:cov}
Let $A_{rq}=\langle Q_r-I,Q_q-I\rangle_{\cH^*}$.
In decreasing block-size order,
\begin{equation}\label{eq:cov}
 A_{rr}=d_{m_r},\quad
 A_{rq}=c_{m_r}:=2d_{m_r+1}-d_{m_r}\quad(r<q),\quad A_{qr}=A_{rq}.
\end{equation}
Thus $e(Q_{\btheta})^2=\btheta^\top A\btheta$.
Every off-diagonal entry in a given upper row is the same.
\end{proposition}
\begin{proof}
Each block is a shift of $H_{m_r}$. For $r<q$, averaging the kernel over
$H_{m_r}$ and $H_{m_q}$ reduces to an average over $H_{m_r}$,
since $H_{m_q}\subset H_{m_r}$.

There is a further binary simplification. The start of any later block
lies in the other half of the same aligned $2N_r$ interval.
Radical inversion gives
\[
 \phi_2(a_q)-\phi_2(a_r)
 \equiv\frac1{2N_r}\pmod{1/N_r}.
\]
The higher binary digits are unchanged, the digit at level $m_r$
changes from zero to one, and the lower digits contribute multiples
of $1/N_r$. Therefore the cross covariance is
\[
 \frac1{N_r}\sum_{j=0}^{N_r-1}
 \kappa\!\left(\bh\frac{j+1/2}{N_r}\right)-1.
\]
The full $2N_r$ lattice splits into $H_{m_r}$ and this other coset.
Averaging the two cosets proves the off-diagonal formula.
The diagonal follows from shift invariance.
\end{proof}
This specializes HKKN12's Proposition 2 and must be acknowledged as
such. It gives the relevant structure without a Toeplitz assumption on
the pointwise Gram matrix in van der Corput order.

\section{An explicit positive-weight recursion}
The usual HKKN12 expression is
$A^{-1}\bbone/(\bbone^\top A^{-1}\bbone)$.
Equation~\eqref{eq:cov} allows the same optimum to be found without
an inverse or even forming $A$.

\begin{theorem}[Optimal binary block rule]\label{thm:recursion}
The unique minimum of $\btheta^\top A\btheta$ over real $\btheta$ with
$\bbone^\top\btheta=1$ has every component positive.
Define $E_L=d_{m_L}$. For $r=L-1,\ldots,1$, compute
\begin{align}
 D_r&=d_{m_r}+E_{r+1}-2c_{m_r},\label{eq:denom}\\
 t_r&=\frac{E_{r+1}-c_{m_r}}{D_r},\label{eq:tr}\\
 E_r&=\frac{d_{m_r}E_{r+1}-c_{m_r}^{\,2}}{D_r}.\label{eq:Er}
\end{align}
Recover the masses in a forward pass:
\[
 p=1;\quad \theta_r=p\,t_r,\quad p\leftarrow p(1-t_r)
 \quad(r=1,\ldots,L-1);\quad\theta_L=p.
\]
For $L=1$, take $\theta_1=1$. The optimum squared error is $E_1$.
\end{theorem}
\begin{proof}
Let $R$ be any normalized rule supported on blocks after $r$.
All its nodes lie in one shifted other coset of $H_{m_r}$ inside
$H_{m_r+1}$. For each $\bk\in\mathcal D_{m_r}$ the characters at
these nodes have the same phase. Normalization, even for signed weights,
implies $|R(e^{2\pi\ii\bk^\top\cdot})|=1$.
Consequently
\begin{equation}\label{eq:tail}
 e(R)^2\ge d_{m_r}.
\end{equation}
Also $\langle Q_r-I,R-I\rangle=c_{m_r}$, independently of its tail masses.

If the best normalized tail has squared error $v=E_{r+1}$, then the
best error with mass $t$ on block $r$ is
\[
 t^2d_{m_r}+2t(1-t)c_{m_r}+(1-t)^2v.
\]
Since $v\ge d_{m_r}>c_{m_r}$, its quadratic coefficient $D_r$ is
positive, and its minimizer is~\eqref{eq:tr}, strictly between zero and
one. Substitution gives~\eqref{eq:Er}. A tail of total mass zero has
zero cross term with $Q_r$ and nonnegative squared norm; it cannot
improve on $d_{m_r}$ and cannot invalidate this interior optimum.
Backward induction proves the recursion and positivity.

For uniqueness, $A$ is positive definite. A vanishing nonzero combination
of block error representers would give a finite atomic signed measure
with all nonconstant Fourier coefficients zero. It would equal a
multiple of Lebesgue measure, which is impossible for a nonzero finite
atomic measure. Strictly positive $\rho_{\bk}$ and distinct nodes
justify this argument.
\end{proof}

\begin{corollary}[Precise gain over truncation]\label{cor:gain}
If $n$ is non-dyadic, so $L\ge2$, and $M=m_1$, then
\begin{equation}\label{eq:gain}
 e(Q_{\rm block})^2
 =d_M-\frac{4(d_M-d_{M+1})^2}
 {d_M+E_2-2(2d_{M+1}-d_M)}
 <d_M=e(Q_{\rm trunc})^2.
\end{equation}
All $n$ sample weights are positive. Also
\[
 e(Q_{\rm full\ opt})\le e(Q_{\rm block})
 \le e(Q_{\rm HKKN}(a)),\qquad
 e(Q_{\rm block})\le e(Q_{\rm equal}).
\]
If $e(Q_{2^m})\le C\,2^{-\beta m}$ for all relevant $m$, then
$e(Q_{\rm block})\le C\,2^\beta n^{-\beta}$.
\end{corollary}
\begin{proof}
Subtract~\eqref{eq:Er} from $d_M$ and use
$d_M-c_M=2(d_M-d_{M+1})$.
The comparisons follow from feasible weight classes.
For the rate use $2^M>n/2$.
\end{proof}
The rate is already available by truncation and HKKN12.
The useful statements here are positive use of all points, a
\emph{quantified} strict gain, and inexpensive computation.
Worst-case-error inequalities do not order actual errors for every
individual integrand. Nor do they imply that unrestricted nodewise
optimal weights are positive.

\subsection{Cost and interrupted computation}
The backward recursion and forward pass cost $O(L)$ arithmetic and
$O(L)$ storage, where $L\le M+1=O(\log n)$.
The values $d_0,\ldots,d_{M+1}$ depend on the kernel and generator,
not on $f$ or the final binary digits of $n$.
Direct use of~\eqref{eq:dm} costs
\[
 O\!\left(s\sum_{m=0}^{M+1}2^m\right)=O(sn),
\]
assuming one kernel evaluation costs $O(s)$.
This includes kernel evaluations on the \emph{next} full lattice;
no extra integrand evaluations are required there.
Cached values give $O(\log n)$ postprocessing at a requested $n$.

Keep one function-value sum per active binary block, merging adjacent
equal-sized blocks as new completed values arrive.
Aggregate storage is $O(\log n)$, and a requested estimate takes
$O(\log n)$ once the weights are known.
Reporting after every point costs $O(n\log n)$ overall, apart from
precomputation and integrand evaluation.
These are arithmetic counts, not a stability theorem.
The kernel-average-minus-one formula for $d_m$ and small differences
may need higher precision or a positive spectral calculation.

The data must form a completed sequence prefix; a failure that loses
values or leaves an arbitrary asynchronous subset is a different
problem. The final prefix length need not be known in advance.

\section{What this accomplishes, and what it does not}
\subsection{A thirteen-point example}
In one dimension with $h=1$ and Hilbert smoothness $\alpha=2$,
\[
 d_m=C\,2^{-4m},\qquad C=2\gamma^4\zeta(4).
\]
For $13=8+4+1$, the optimal block masses are
\[
 (\theta_1,\theta_2,\theta_3)
 =\left(\frac{613}{682},\frac{137}{1364},\frac1{1364}\right).
\]
Each of the first eight points has weight $613/5456$,
each of the next four has weight $137/5456$,
and the last point has weight $4/5456$. Exactly,
\[
 e(Q_{\rm block})^2=C\,\frac{4421}{22347776},\qquad
 \frac{e(Q_{\rm block})}{e(Q_{\rm trunc})}
 =\sqrt{\frac{4421}{5456}}\approx0.9002.
\]
This is a sanity check, not high-dimensional evidence.
Positive weights on surplus points can be very small; report that
rather than equating positive use with a large practical benefit.

\subsection{Do not promise a universal factor-free bound}
For the same one-dimensional kernel at $n=3$, optimal block masses
are $(9/10,1/10)$ and point weights are $(9/20,9/20,1/10)$.
The squared error is $13C/256$.
The exact dyadic law has leading constant $\sqrt C$ in
$e(Q_{2^m})=\sqrt C\,(2^m)^{-2}$.
Keeping that \emph{same} constant for $n=3$ would require squared
error at most $C/81$, but $13/256>1/81$.

Thus~\eqref{eq:gain} improves the actual truncated Hilbert error;
it does not remove the entire $(n/2^M)^\alpha$ loss in the strongest
sense. Referee Z's displayed (T) concerns $P_\alpha$ and the quoted
Banach bound. This note addresses its Hilbert analogue; it does not
prove a sharper $P_\alpha$ bound by changing notation.

\subsection{A Banach-space obstruction}
In the one-dimensional Banach space of the manuscript, put
$C_\alpha=2\gamma^\alpha\zeta(\alpha)$ and
$q=2^{-\alpha}<1/2$.
For $n=3$, assign mass $1-t$ to the first two points and $t$ to
the singleton. Split Fourier modes into odd integers, integers
divisible by two but not four, and multiples of four. This gives
\begin{equation}\label{eq:banach}
 P_\alpha(t)=C_\alpha
 [q^2+q(1-q)|1-2t|+(1-q)|t|].
\end{equation}
On $0\le t\le1/2$ its increase from $t=0$ is
$C_\alpha(1-q)(1-2q)t>0$.
For $t<0$ it decreases towards zero; for $t>1/2$ it increases.
The unique minimum over real $t$ is truncation, $t=0$.
This block family cannot strictly improve the exact Banach truncation
error at every non-dyadic $n$, even though it can in the Hilbert space.
That distinction should shape the revised paper.

\section{Experiments that would make the revision persuasive}
This is a proposed protocol, not completed large-scale experimentation.
The companion exact-arithmetic check validates only the small identities
and examples.

\begin{enumerate}
\item \textbf{Methods.}
Use identical nodes and the same kernel for equal point weights;
truncation at $2^{\lfloor\log_2n\rfloor}$; HKKN12 power masses with
$a=2,3,4$; the recursive optimal block rule; and, at smaller $n$,
the full nodewise optimum with weights summing to one.
Also compare the recursion with the dense $L\times L$ HKKN12 solve.
They must agree: these are two computations of the same rule.

\item \textbf{Smoothness, dimensions and budgets.}
Use Hilbert $\alpha=2$ as the primary case, add $\alpha=3$, and retain
$\alpha=1$ only as context.
For example, use $s=2,6,20$ with two declared coordinate-weight decays.
Scan every $n$ to a modest horizon such as $2^{12}$, then prespecify
larger budgets. Include $2^M+1$, $2^{M+1}-1$, $13\cdot2^p$, and
different binary digit counts, rather than favorable prefixes only.
A common finite generator need not be good for every kernel:
report the common-generator comparison and, if available,
tuned-generator comparisons separately.

\item \textbf{Worst-case-error outputs.}
Plot $e_{\rm method}/e_{\rm trunc}$ against both $n$ and
$n/2^{\lfloor\log_2n\rfloor}$, marking powers of two.
Include ratios to HKKN12 power weights and to the full nodewise
optimum. Report ratio distributions and maxima, not only chosen curves.
Separate $e$ and $e^2$. Time precomputation, weight calculation, and
application separately; record memory and arithmetic precision.
Agreement with the old optimum is verification; the gain in cost is
the computational comparison.

\item \textbf{Integrands and uncertainty.}
Use smooth periodic products
$f(\bx)=\prod_j[1+\eta_j\cos(2\pi x_j)]$, with integral one and finite
Fourier support, and a broader-spectrum periodic example with a
checked finite Hilbert norm.
For actual-integrand comparisons use paired common random shifts,
for example at least 32 independent shifts, with absolute-error
summaries. Shifts do not change this periodic worst-case error,
so they are not replicate worst-case-error experiments.
Keister may be an application illustration, but the transformed
integrand's endpoint regularity and norm must be checked before
claiming it belongs to the theoretical space.
Protect actual-error ratios against a zero truncated denominator.

\item \textbf{Reproducibility and numerical checks.}
Record generator entries or an archived file and checksum, residue
horizon, ordering, kernel scaling, coordinate weights, seeds, software,
precision, and scripts.
Check $\sum_iw_i=1$, positivity,~\eqref{eq:cov}, and
$A\btheta-E_1\bbone$.
Validate tiny discrepancies in higher precision; do not silently clip
negative roundoff discrepancies to zero.
\end{enumerate}

\section{Changes needed to address all three reports}
\paragraph{Positioning (Y and Z).}
State truncation and the trivial full-optimum monotonicity at the start.
Treat the rate alone as background. Explain the distinction from HKKN12
at both theorem and algorithm levels, why the Hilbert norm is relevant,
and what practical gain warrants weighted surplus values.
A small improvement in one example does not establish readiness
for resubmission.

\paragraph{Correctness and scope (X and Y).}
Specify $\alpha>1$ where Banach Fourier summability is required.
The replacement of $r(k,\gamma_j)$ by $\max(1,|k|/\gamma_j)$
needs $\gamma_j\le1$; retain that assumption or remove the replacement.
Restore the missing $(b-1)$ in the logarithm-free bound.
State existence and generator quantifiers consistently in inherited
bounds, and audit base restrictions against HN03 \cite{HN03}.

For general base $b$, a repeated digit represents several blocks with
different shifts, not a single phase times a digit.
With $a_\ell=\sum_{t>\ell}n_tb^t$, the correct character decomposition is
\[
 \frac{e^{2\pi\ii\bk^\top\bDelta}}{n}
 \sum_{\ell=0}^M\sum_{q=0}^{n_\ell-1}
 b^\ell\bbone_{\{\bk^\top\bh=0\bmod b^\ell\}}
 e^{2\pi\ii\bk^\top\bh\phi_b(a_\ell+qb^\ell)}.
\]
This preserves the triangle bound and covers Y's $b=3,n=2$ example.
Correct the misplaced power in the denominator, extra parentheses,
duplicated alignment end, and missing wavenumber inner product
in the Hilbert Fourier expansion.

\paragraph{Normalization and the visible numerical code.}
For unnormalized weights the identity is
$e(w)^2=1-2\bbone^\top w+w^\top Kw$.
Only with $\bbone^\top w=1$ may it become $w^\top Kw-1$.
The visible \texttt{Programs/discrepancy.m} currently solves
$u=K^{-1}\bbone$ and uses the former expression without normalization.
That differs from the constrained optimum derived in the manuscript.
Audit figure provenance before interpreting existing plots.
The constrained optimum uses $w=u/(\bbone^\top u)$;
solve a system rather than explicitly forming an inverse.

Also, \texttt{GramMat.m} uses a $B_2$ factor with coefficient
$\texttt{coordwts}_j/2$.
Under the manuscript's $\alpha=1$ convention the coefficient is
$2\pi^2\gamma_j^2$. They agree only when
$\texttt{coordwts}_j=4\pi^2\gamma_j^2$.
For integer Hilbert $\alpha$, the intended factor is
\[
 1+\frac{(-1)^{\alpha+1}(2\pi\gamma_j)^{2\alpha}}{(2\alpha)!}
 B_{2\alpha}(\{x_j-y_j\}).
\]
The existing $B_2$ routine does not supply an $\alpha=2$ experiment.

\paragraph{Presentation and reproducibility (especially X).}
Describe rank-1 rules as a class of lattice rules, rather than their
historical definition.
Explain generalized integers through coherent residues, and distinguish
one-dimensional radical inversion from its multidimensional image.
Define the dual set before using it.
Use a declared $\ii^2=-1$ and consistent arguments and brackets.
H\"older alone gives an inequality; retain a worst-case equality only
after a Fourier-phase attainability argument. For $\alpha>1$ the
candidate coefficients are absolutely summable, so this can be done.
Replace the unsupported assertion that no finite formula exists by a
statement about the unavailable formula used here.

Give actual generators, shifts, dimensions, and plotted quantities.
Improve Figure 1's fonts and contrast; make captions consistent.
Distinguish existence from construction, stronger norm from larger
unit ball, exact formula from bound, and illustration from general
evidence. End experiments with supported conclusions.
Remove the false generic Toeplitz claim and the cost derived from it:
a dense nodewise calculation costs $O(sn^2+n^3)$;
a full dyadic lattice has circulant structure after permutation.
The block recursion has a different, explicitly justified cost.

\section{Suggested manuscript structure and decision point}
A possible title is \emph{Positive Block Weights for Lattice Sequence
Prefixes: Exact Error Reduction and Fast Computation}.
An outline is:
\begin{enumerate}
\item Arbitrary stopping, the truncation baseline, and the specific
gap beyond existing compound rules.
\item Binary blocks and the periodic Hilbert setting, with HN03 and
HKKN12 positioned explicitly.
\item Covariance identity, positive recursion, exact gain over
truncation, limitations, and numerical stability.
\item Reproducible $\alpha>1$ comparisons with old compound rules
and full nodewise weights.
\item Practical consequences and open Banach-space questions.
Move long equal-weight bounds to an appendix if useful.
\end{enumerate}
Megan's current work could connect to the covariance calculation,
stable evaluation of $d_m$, and the $\alpha>1$ benchmarks.
This is a proposed connection, not an assumption about her results,
an assignment, or an authorship decision.

Before choosing a resubmission venue, jointly verify the proofs,
check whether the specialization and positivity result are already
known, and establish whether computational and empirical gains are
substantial. If this is only a short specialization with small gains,
a shorter paper or a more extensive application study may be more
convincing than retaining the original scope.

\begin{thebibliography}{9}
\bibitem{HN03}
F.\ J.\ Hickernell and H.\ Niederreiter,
The existence of good extensible rank-1 lattices,
\emph{Journal of Complexity} \textbf{19} (2003), 286--300.
\bibitem{HKKN12}
F.\ J.\ Hickernell, P.\ Kritzer, F.\ Y.\ Kuo and D.\ Nuyens,
Weighted compound integration rules with higher order convergence
for all $N$,
\emph{Numerical Algorithms} \textbf{59} (2012), 161--183.
\href{https://doi.org/10.1007/s11075-011-9482-5}{doi:10.1007/s11075-011-9482-5}.
\end{thebibliography}
\paragraph{Sources and validation status.}
Prepared from the current manuscript, all three reports in
\texttt{Referee\_Reports}, and the published article supplied as
\begin{quote}\small\texttt{relevant articles/WeightedCompoundIntegration.pdf}\end{quote}
The proofs are drafts for author review.
The companion \texttt{verify\_binary\_blocks.py} checks small
Bernoulli-kernel cases in exact rational arithmetic; it does not
replace proofs, a novelty search, or the proposed numerical study.
\end{document}
